%%%%%%%%%%%%%%%%%%%%%%%%%%%%%%%%%%%%%%%%%%%%%%%%%%%%%%%%%%%%%%%%%%%%%%%%%%%%%%%
% APPENDIX W: Akerlof, Spence & Stiglitz – Information Economics
% The Informational Foundations of Context-Dependent Markets
%%%%%%%%%%%%%%%%%%%%%%%%%%%%%%%%%%%%%%%%%%%%%%%%%%%%%%%%%%%%%%%%%%%%%%%%%%%%%%%
% Category: LIT
% Language: English
%%%%%%%%%%%%%%%%%%%%%%%%%%%%%%%%%%%%%%%%%%%%%%%%%%%%%%%%%%%%%%%%%%%%%%%%%%%%%%%

% =============================================================================
\begin{tcolorbox}[colback=blue!5!white, colframe=blue!75!black,
    title=Quick Reference: Appendix UNMAPPED_INF]
\textbf{Appendix:} INF (LIT)

\textbf{Core Question:} What are the key research findings in this literature domain?

\textbf{Key Concepts:}
\begin{itemize}[nosep]
\item Literature synthesis and integration
\item Framework application and validation
\item Cross-domain connections
\end{itemize}

\textbf{Cross-References:} See Chapter linkage below
\end{tcolorbox}


\begin{abstract}
This appendix provides a systematic review and integration of key research findings
in the LIT domain. It synthesizes literature relevant to the Evidence-Based
Framework (EBF) and identifies connections to the 10C CORE architecture.
\end{abstract}


\begin{tcolorbox}[
    colback=orange!5!white,
    colframe=orange!75!black,
    title={\textbf{Appendix Scope: Ziel / In-Scope / Out-of-Scope / Constraints}},
    fonttitle=\bfseries
]

\textbf{Ziel (Objective):}
\begin{itemize}[nosep]
    \item Synthesize key research findings for LIT integration
\end{itemize}

\textbf{In-Scope:}
\begin{itemize}[nosep]
    \item Literature review and synthesis
    \item Parameter estimates from empirical studies
    \item Framework integration points
\end{itemize}

\textbf{Out-of-Scope:}
\begin{itemize}[nosep]
    \item Original empirical analysis $\rightarrow$ METHOD appendices
    \item Detailed mathematical derivations $\rightarrow$ FORMAL appendices
\end{itemize}

\textbf{Constraints:}
\begin{itemize}[nosep]
    \item Literature selection based on citation impact and relevance
    \item Parameter estimates subject to study conditions
\end{itemize}

\end{tcolorbox}

\section{The Fundamental Question}
\label{sec:inf-fundamental}
%%%%%%%%%%%%%%%%%%%%%%%%%%%%%%%%%%%%%%%%%%%%%%%%%%%%%%%%%%%%%%%%%%%%%%%%%%%%%%%

\begin{tcolorbox}[
    colback=red!5!white,
    colframe=red!75!black,
    title={\textbf{Fundamental Question}}
]
\textbf{How does unknown integrate with and contribute to the
Evidence-Based Framework for understanding behavioral outcomes?}

This appendix addresses the core challenge of formalizing behavioral principles
within a rigorous theoretical framework while maintaining practical applicability.
\end{tcolorbox}


\section{Akerlof, Spence \& Stiglitz: The Informational Foundations}
\label{app:information-economics}

%==============================================================================
% PART 0: INTEGRATION OVERVIEW
%==============================================================================
\subsection{Framework Integration Map}

This appendix demonstrates how the research program of George Akerlof, Michael Spence, 
and Joseph Stiglitz provides the \textbf{informational foundations} for the 
$(C, \Psi, C^*, K, Q)$ framework. Their work establishes that information asymmetry 
fundamentally alters market outcomes and that $\Psi_K$ (the informational dimension 
of context) is not merely a parameter but a determinant of whether coherent 
configurations $C^*$ can exist at all.

\begin{center}
\renewcommand{\arraystretch}{1.4}
\begin{tabular}{|l|l|l|}
\hline
\textbf{Framework Element} & \textbf{A/S/S Contribution} & \textbf{Key Papers} \\
\hline
$\Psi_K$ (Informational Context) & Information as primitive & All foundational papers \\
$C^*$ Existence & Markets can fail to clear & Akerlof 1970, Stiglitz-Weiss 1981 \\
$C_{ij}^{asym} \neq C_{ij}^{sym}$ & Asymmetry changes structure & Spence 1973, Rothschild-Stiglitz 1976 \\
Endogenous $\Psi_K$ & Signaling/Screening & Spence 1973, Stiglitz 1975 \\
$K \to 0$ Possible & Market unraveling & Akerlof 1970, Wilson 1980 \\
Policy under $\Psi_K$ heterogeneity & Intervention design & Stiglitz 2002, Greenwald-Stiglitz 1986 \\
\hline
\end{tabular}
\end{center}

\paragraph{The Information Trio's Unique Role.}
While other contributors address specific framework components (Kahneman/Tversky: 
psychological $\Psi_C$; Bloom: uncertainty $\Psi_T$; Acemoglu: institutional $\Psi_I$), 
Akerlof, Spence, and Stiglitz uniquely establish:
\begin{enumerate}
    \item \textbf{Information asymmetry} as market-defining context
    \item \textbf{Market failure} as equilibrium outcome (not just inefficiency)
    \item \textbf{Endogenous information} via signaling and screening
    \item \textbf{Multiple equilibria} depending on $\Psi_K$ configuration
    \item \textbf{Policy implications} when first-best is unattainable
\end{enumerate}

%==============================================================================
% PART I: ADVERSE SELECTION – WHEN C* CANNOT EXIST
%==============================================================================
\subsection{Part I: Adverse Selection – When $C^*$ Cannot Exist}

\subsubsection{Theoretical Foundation}

Akerlof's (1970) fundamental insight: when quality is unobservable, 
markets can unravel entirely.

\begin{equation}
\boxed{
\Psi_K^{low} \Rightarrow C^* = \emptyset \quad \text{(no equilibrium with trade)}
}
\label{eq:akerlof-unraveling}
\end{equation}

\paragraph{The Mechanism.}
\begin{enumerate}
    \item Sellers know quality; buyers don't ($\Psi_K$ asymmetric)
    \item Buyers offer average-quality price
    \item High-quality sellers exit (price too low)
    \item Average quality falls; price falls
    \item Iteration $\to$ only lowest quality remains or market collapses
\end{enumerate}

\subsubsection{Framework Translation}

In standard Arrow-Debreu, $C^*$ always exists (complete markets, full information).
Akerlof shows:
\begin{equation}
C^*_{AD} \text{ exists} \quad \text{but} \quad C^*_{Akerlof}(\Psi_K^{asym}) \text{ may not exist}
\end{equation}

This is not inefficiency within equilibrium—it is \textbf{equilibrium nonexistence}.

\subsubsection{The Lemons Model Formally}

Let $q \in [0,1]$ be quality, uniformly distributed. Seller valuation: $v_s(q) = q$.
Buyer valuation: $v_b(q) = \frac{3}{2}q$. Efficient outcome: all cars trade.

With asymmetric information:
\begin{align}
\text{Buyer offers price } p &\Rightarrow \text{Sellers with } q \leq p \text{ accept} \\
\text{Average quality } \bar{q}(p) &= \frac{p}{2} \\
\text{Buyer's expected value } &= \frac{3}{2} \cdot \frac{p}{2} = \frac{3p}{4} < p
\end{align}

No price $p > 0$ is sustainable $\Rightarrow$ $C^* = \emptyset$.

\subsubsection{$\Psi_K$ as Existence Condition}

The framework implication:
\begin{equation}
\exists C^* \quad \Leftrightarrow \quad \Psi_K > \Psi_K^{threshold}
\end{equation}

Market viability depends on informational context. This explains why:
\begin{itemize}
    \item Used car markets require warranties, inspections, reputation (raising $\Psi_K$)
    \item Insurance markets require risk classification (reducing asymmetry)
    \item Credit markets require collateral, credit scores (signaling/screening)
\end{itemize}

%==============================================================================
% PART II: SIGNALING – ENDOGENOUS Ψ_K
%==============================================================================
\subsection{Part II: Signaling – Endogenous $\Psi_K$}

\subsubsection{Theoretical Foundation}

Spence (1973) shows that informed parties can \textit{reveal} information through 
costly actions:

\begin{equation}
\boxed{
\Psi_K^{post-signal} > \Psi_K^{pre-signal} \quad \text{via costly signaling}
}
\label{eq:spence-signaling}
\end{equation}

\paragraph{Key Insight.}
Information is not exogenous. Agents invest in $\Psi_K$ improvement when:
\begin{equation}
\text{Benefit of revelation} > \text{Cost of signal}
\end{equation}

\subsubsection{The Education Signaling Model}

Workers have ability $\theta \in \{\theta_L, \theta_H\}$. Education $e$ costs $c(\theta, e)$
with $\frac{\partial c}{\partial \theta} < 0$ (cheaper for high ability).

\textbf{Separating Equilibrium:}
\begin{align}
e^*_H &: \quad w(e^*_H) - c(\theta_H, e^*_H) \geq w(0) \\
e^*_L &= 0 : \quad w(0) \geq w(e^*_H) - c(\theta_L, e^*_H)
\end{align}

High types signal; low types don't $\Rightarrow$ $\Psi_K$ increases.

\subsubsection{Framework Translation}

Signaling makes $\Psi_K$ endogenous:
\begin{equation}
\Psi_K = \Psi_K^{exogenous} + f(\text{signaling investment})
\end{equation}

This has profound implications:
\begin{itemize}
    \item $\Psi_K$ is not just context—it's a choice variable
    \item Equilibrium $\Psi_K$ depends on cost structures
    \item Multiple equilibria possible (pooling vs. separating)
\end{itemize}

\subsubsection{Signaling Beyond Education}

\begin{center}
\renewcommand{\arraystretch}{1.3}
\begin{tabular}{|l|l|l|}
\hline
\textbf{Domain} & \textbf{Signal} & \textbf{$\Psi_K$ Revealed} \\
\hline
Labor & Education, credentials & Ability, commitment \\
Finance & Dividends, debt & Firm quality \\
Products & Warranties, advertising & Product quality \\
Social & Conspicuous consumption & Wealth, status \\
\hline
\end{tabular}
\end{center}

%==============================================================================
% PART III: SCREENING – Ψ_K EXTRACTION BY UNINFORMED
%==============================================================================
\subsection{Part III: Screening – $\Psi_K$ Extraction by Uninformed}

\subsubsection{Theoretical Foundation}

Rothschild \& Stiglitz (1976) show the uninformed party can design menus 
that induce self-selection:

\begin{equation}
\boxed{
\text{Menu } \{(coverage_i, premium_i)\} \Rightarrow \text{Type } i \text{ self-selects}
}
\label{eq:rs-screening}
\end{equation}

\subsubsection{The Insurance Screening Model}

Insurers offer contracts. High-risk types ($H$) and low-risk types ($L$) have 
different accident probabilities $p_H > p_L$.

\textbf{Separating Equilibrium:}
\begin{itemize}
    \item $H$-types get full insurance at actuarially fair price $p_H$
    \item $L$-types get partial insurance at price $p_L$ (distorted to prevent $H$ mimicking)
\end{itemize}

\textbf{Key Result:} Pooling equilibrium cannot exist—it would be ``cream-skimmed.''

\subsubsection{Framework Translation}

Screening extracts $\Psi_K$ through mechanism design:
\begin{equation}
\Psi_K^{extracted} = g(\text{menu complexity}, \text{type heterogeneity})
\end{equation}

\paragraph{Critical Insight: Equilibrium May Not Exist.}
Rothschild-Stiglitz showed that with continuous type distributions, 
\textit{no equilibrium may exist}—a fundamental challenge to standard theory.

\subsubsection{Screening Applications}

\begin{center}
\renewcommand{\arraystretch}{1.3}
\begin{tabular}{|l|l|l|}
\hline
\textbf{Market} & \textbf{Screening Device} & \textbf{Information Extracted} \\
\hline
Insurance & Deductibles, coverage limits & Risk type \\
Credit & Interest rates, collateral & Default probability \\
Airlines & Fare classes, restrictions & Price sensitivity \\
Telecom & Data plans, contract length & Usage patterns \\
\hline
\end{tabular}
\end{center}

%==============================================================================
% PART IV: CREDIT RATIONING – C* WITH EXCESS DEMAND
%==============================================================================
\subsection{Part IV: Credit Rationing – $C^*$ with Excess Demand}

\subsubsection{Theoretical Foundation}

Stiglitz \& Weiss (1981) demonstrate that credit markets may not clear even 
in equilibrium:

\begin{equation}
\boxed{
r^* < r^{market-clearing} \quad \text{and} \quad D(r^*) > S(r^*) \quad \text{in equilibrium}
}
\label{eq:sw-rationing}
\end{equation}

\subsubsection{The Mechanism}

Higher interest rates have two effects:
\begin{enumerate}
    \item \textbf{Direct:} Revenue per loan increases
    \item \textbf{Adverse selection:} Safe borrowers exit; risky borrowers remain
    \item \textbf{Moral hazard:} Borrowers take more risk
\end{enumerate}

Above some $r^*$, the negative selection/incentive effects dominate:
\begin{equation}
\frac{\partial \pi_{bank}}{\partial r} < 0 \quad \text{for } r > r^*
\end{equation}

\subsubsection{Framework Translation}

Credit rationing shows $C^*$ can involve \textbf{non-price allocation}:
\begin{equation}
C^*_{credit} = (r^*, \text{rationing rule}) \neq (r^{clearing}, \text{no rationing})
\end{equation}

This fundamentally challenges the price-mechanism view of markets.

\paragraph{Implications for $\Psi_K$.}
\begin{itemize}
    \item Information ($\Psi_K$) determines rationing severity
    \item Better credit scoring reduces rationing (raises $\Psi_K$)
    \item Collateral requirements are screening devices
    \item Relationship banking is $\Psi_K$ investment
\end{itemize}

%==============================================================================
% PART V: MARKET FAILURE AS EQUILIBRIUM
%==============================================================================
\subsection{Part V: Greenwald-Stiglitz Theorem – Pervasive Market Failure}

\subsubsection{Theoretical Foundation}

Greenwald \& Stiglitz (1986) prove that with asymmetric information, 
markets are \textit{generically} inefficient:

\begin{equation}
\boxed{
\Psi_K < \Psi_K^{complete} \Rightarrow C^* \text{ is constrained Pareto inefficient}
}
\label{eq:gs-theorem}
\end{equation}

\subsubsection{The Theorem}

\textbf{Greenwald-Stiglitz (1986):} Whenever information is imperfect or markets 
incomplete, the competitive equilibrium is not constrained Pareto efficient. 
There exist government interventions (taxes, subsidies) that can make everyone 
better off.

\paragraph{Implication.}
The ``invisible hand'' fails not occasionally but \textit{generically} under 
realistic information conditions.

\subsubsection{Framework Translation}

The theorem establishes:
\begin{equation}
Q(C^*(\Psi_K^{real})) < Q(C^*(\Psi_K^{complete})) \quad \text{generically}
\end{equation}

Policy can improve outcomes even without resolving information asymmetry—by 
correcting the externalities that asymmetry creates.

%==============================================================================
% PART VI: SYNTHESIS – THE INFORMATION ECONOMICS EQUATIONS
%==============================================================================
\subsection{Part VI: Synthesis – The Information Economics Equations}

\subsubsection{Complete Framework Translation}

\begin{center}
\renewcommand{\arraystretch}{1.5}
\begin{tabular}{|l|c|l|}
\hline
\textbf{A/S/S Finding} & \textbf{Equation} & \textbf{Framework Element} \\
\hline
Adverse selection & $\Psi_K^{low} \Rightarrow C^* = \emptyset$ & Existence depends on $\Psi_K$ \\
\hline
Signaling & $\Psi_K = f(\text{costly actions})$ & $\Psi_K$ is endogenous \\
\hline
Screening & $\Psi_K^{extracted} = g(\text{menu})$ & Mechanisms reveal information \\
\hline
Credit rationing & $C^* \neq$ market-clearing & Non-price allocation in equilibrium \\
\hline
Greenwald-Stiglitz & $Q(C^*) < Q^{first-best}$ & Generic inefficiency \\
\hline
Multiple equilibria & $C^*_1 \neq C^*_2$ both stable & $\Psi_K$ determines selection \\
\hline
Asymmetric $C_{ij}$ & $C_{ij}^{asym} \neq C_{ij}^{sym}$ & Information changes complementarities \\
\hline
\end{tabular}
\end{center}

\subsubsection{The Information Trio's Unique Contribution}

\textbf{Without Akerlof/Spence/Stiglitz, the framework would:}
\begin{itemize}
    \item Assume $C^*$ always exists (Arrow-Debreu)
    \item Treat $\Psi_K$ as exogenous parameter
    \item Miss credit rationing and non-price allocation
    \item Lack foundation for market failure as \textit{equilibrium}
    \item Have no theory of signaling/screening dynamics
\end{itemize}

\textbf{With Akerlof/Spence/Stiglitz, the framework gains:}
\begin{itemize}
    \item $\Psi_K$ as existence condition for $C^*$
    \item Endogenous information through signaling/screening
    \item Multiple equilibria depending on $\Psi_K$
    \item Theoretical foundation for policy intervention
    \item Understanding of market unraveling ($K \to 0$)
\end{itemize}

%==============================================================================
% PART VII: COMPLETE PAPER DOCUMENTATION
%==============================================================================
\subsection{Part VII: Complete Paper Documentation}

%--- AKERLOF ---
\subsubsection{George Akerlof: Adverse Selection and Identity}

\paragraph{Paper 1: Akerlof (1970)}

``The Market for `Lemons': Quality Uncertainty and the Market Mechanism.''
\textit{Quarterly Journal of Economics}, 84(3), 488--500.

\textbf{Summary.}
The foundational paper on adverse selection. Shows how information asymmetry 
about quality can cause market collapse. Used car market as paradigm case.

\textbf{Framework Translation.}
\begin{equation}
\Psi_K^{quality} = 0 \Rightarrow C^*_{trade} = \emptyset
\end{equation}
Market existence requires minimum information quality.

\textbf{Key Finding.}
Nobel Prize paper. Rejected by AER and RES before QJE acceptance. 
Most influential paper in information economics.

\paragraph{Paper 2: Akerlof (1976)}

``The Economics of Caste and of the Rat Race and Other Woeful Tales.''
\textit{Quarterly Journal of Economics}, 90(4), 599--617.

\textbf{Summary.}
Extends adverse selection logic to labor markets, showing how caste systems 
and inefficient signaling equilibria can persist.

\textbf{Framework Translation.}
Multiple $C^*$ exist; society can be trapped in inferior equilibrium:
\begin{equation}
C^*_{caste} \prec C^*_{merit} \quad \text{but both stable}
\end{equation}

\paragraph{Paper 3: Akerlof (1980)}

``A Theory of Social Custom, of Which Unemployment May Be One Consequence.''
\textit{Quarterly Journal of Economics}, 94(4), 749--775.

\textbf{Summary.}
Social norms as equilibria. Disobedience is costly only because others 
conform—self-fulfilling beliefs create multiple equilibria.

\textbf{Framework Translation.}
$\Psi_S$ (social context) creates coordination on $C^*$:
\begin{equation}
C^* = C^*(\Psi_S) \quad \text{where } \Psi_S \text{ includes norms}
\end{equation}

\paragraph{Paper 4: Akerlof \& Kranton (2000)}

``Economics and Identity.''
\textit{Quarterly Journal of Economics}, 115(3), 715--753.

\textbf{Summary.}
Introduces identity into utility function. People derive utility from 
conforming to identity-group norms.

\textbf{Framework Translation.}
Identity ($U^{IDN}$) is FEPSDE component:
\begin{equation}
U = U^F + U^E + U^P + U^S + U^D + U^{IDN}
\end{equation}

\paragraph{Paper 5: Akerlof \& Yellen (1990)}

``The Fair Wage-Effort Hypothesis and Unemployment.''
\textit{Quarterly Journal of Economics}, 105(2), 255--283.

\textbf{Summary.}
Workers reduce effort if paid below ``fair'' wage. Explains wage rigidity 
and unemployment through fairness concerns.

\textbf{Framework Translation.}
Reference dependence in wages:
\begin{equation}
e = e(w - w^{fair}) \quad \text{where } w^{fair} = r(\Psi_S)
\end{equation}

\paragraph{Paper 6: Akerlof \& Shiller (2009)}

\textit{Animal Spirits: How Human Psychology Drives the Economy.}
Princeton University Press.

\textbf{Summary.}
Argues macroeconomics needs psychological realism: confidence, fairness, 
corruption, money illusion, stories.

\textbf{Framework Translation.}
$\Psi$ includes psychological factors that drive aggregate outcomes:
\begin{equation}
Y = Y(\Psi_{confidence}, \Psi_{fairness}, \Psi_{narrative}, \ldots)
\end{equation}

%--- SPENCE ---
\subsubsection{Michael Spence: Signaling Theory}

\paragraph{Paper 7: Spence (1973)}

``Job Market Signaling.''
\textit{Quarterly Journal of Economics}, 87(3), 355--374.

\textbf{Summary.}
The foundational signaling paper. Education may be valued not for human 
capital but for sorting workers by unobservable ability.

\textbf{Framework Translation.}
\begin{equation}
\Psi_K^{post-education} > \Psi_K^{pre-education} \quad \text{via costly signal}
\end{equation}
Information is produced, not just transmitted.

\textbf{Key Finding.}
Nobel Prize paper. Created entire field of signaling theory.

\paragraph{Paper 8: Spence (1974)}

\textit{Market Signaling: Informational Transfer in Hiring and Related 
Screening Processes.}
Harvard University Press.

\textbf{Summary.}
Book-length treatment of signaling theory with extensions to warranties, 
advertising, and other domains.

\textbf{Framework Translation.}
Signals work across domains:
\begin{equation}
\Psi_K = \sum_i \alpha_i \cdot \text{Signal}_i
\end{equation}

\paragraph{Paper 9: Spence (1976)}

``Competition in Salaries, Credentials, and Signaling Prerequisites for Jobs.''
\textit{Quarterly Journal of Economics}, 90(1), 51--74.

\textbf{Summary.}
Credential competition can be socially wasteful. Signaling equilibrium 
may be Pareto dominated by pooling.

\textbf{Framework Translation.}
$C^*_{signaling}$ may have lower $Q$ than $C^*_{pooling}$:
\begin{equation}
Q(C^*_{separating}) < Q(C^*_{pooling}) \quad \text{possible}
\end{equation}
But $C^*_{pooling}$ may not be stable.

\paragraph{Paper 10: Spence (2002)}

``Signaling in Retrospect and the Informational Structure of Markets.''
\textit{American Economic Review}, 92(3), 434--459.

\textbf{Summary.}
Nobel lecture. Reflects on signaling theory's development and applications 
across economics.

\textbf{Framework Translation.}
Synthesis of how signaling affects $\Psi_K$ across market contexts.

%--- STIGLITZ ---
\subsubsection{Joseph Stiglitz: Screening, Rationing, and Policy}

\paragraph{Paper 11: Rothschild \& Stiglitz (1976)}

``Equilibrium in Competitive Insurance Markets: An Essay on the Economics 
of Imperfect Information.''
\textit{Quarterly Journal of Economics}, 90(4), 629--649.

\textbf{Summary.}
Shows competitive insurance markets may have no equilibrium with heterogeneous 
risk types. Screening via self-selection.

\textbf{Framework Translation.}
\begin{equation}
\exists C^* \quad \text{not guaranteed under } \Psi_K^{asymmetric}
\end{equation}

\textbf{Key Finding.}
Equilibrium nonexistence result. Foundational for insurance economics.

\paragraph{Paper 12: Stiglitz \& Weiss (1981)}

``Credit Rationing in Markets with Imperfect Information.''
\textit{American Economic Review}, 71(3), 393--410.

\textbf{Summary.}
Banks may ration credit even with excess demand. Interest rate increases 
cause adverse selection and moral hazard.

\textbf{Framework Translation.}
\begin{equation}
C^*_{credit} \neq \text{market-clearing} \quad \text{under } \Psi_K^{borrower} > \Psi_K^{lender}
\end{equation}

\textbf{Key Finding.}
Explains credit crunches, financial exclusion. 8,000+ citations.

\paragraph{Paper 13: Greenwald \& Stiglitz (1986)}

``Externalities in Economies with Imperfect Information and Incomplete Markets.''
\textit{Quarterly Journal of Economics}, 101(2), 229--264.

\textbf{Summary.}
Proves that markets with asymmetric information are generically constrained 
Pareto inefficient.

\textbf{Framework Translation.}
\begin{equation}
\Psi_K < \Psi_K^{complete} \Rightarrow Q(C^*) < Q^{first-best} \quad \text{generically}
\end{equation}

\textbf{Key Finding.}
The ``Greenwald-Stiglitz Theorem.'' Provides foundation for policy intervention.

\paragraph{Paper 14: Stiglitz (1975)}

``The Theory of `Screening,' Education, and the Distribution of Income.''
\textit{American Economic Review}, 65(3), 283--300.

\textbf{Summary.}
Distinguishes screening from signaling. Uninformed parties design tests 
to extract information.

\textbf{Framework Translation.}
\begin{equation}
\Psi_K^{uninformed} \to \Psi_K^{informed} \quad \text{via screening design}
\end{equation}

\paragraph{Paper 15: Stiglitz (1987)}

``The Causes and Consequences of the Dependence of Quality on Price.''
\textit{Journal of Economic Literature}, 25(1), 1--48.

\textbf{Summary.}
Comprehensive survey of why price can signal quality. Links adverse selection, 
moral hazard, and efficiency wages.

\textbf{Framework Translation.}
Price is informative:
\begin{equation}
\mathbb{E}[q | p] = h(p) \quad \text{where } h'(p) > 0 \text{ possible}
\end{equation}

\paragraph{Paper 16: Stiglitz \& Weiss (1983)}

``Incentive Effects of Terminations: Applications to the Credit and Labor Markets.''
\textit{American Economic Review}, 73(5), 912--927.

\textbf{Summary.}
Threat of termination can substitute for monitoring. Unemployment and 
credit denial serve incentive functions.

\textbf{Framework Translation.}
Non-market mechanisms support $C^*$:
\begin{equation}
K_{employment} > K_{spot-market} \quad \text{due to termination threat}
\end{equation}

\paragraph{Paper 17: Shapiro \& Stiglitz (1984)}

``Equilibrium Unemployment as a Worker Discipline Device.''
\textit{American Economic Review}, 74(3), 433--444.

\textbf{Summary.}
Efficiency wage model. Firms pay above market-clearing to deter shirking; 
involuntary unemployment results.

\textbf{Framework Translation.}
\begin{equation}
C^*_{labor} = (w^* > w^{clearing}, U > 0)
\end{equation}
Unemployment is equilibrium feature, not disequilibrium.

\paragraph{Paper 18: Stiglitz (2000)}

``The Contributions of the Economics of Information to Twentieth Century Economics.''
\textit{Quarterly Journal of Economics}, 115(4), 1441--1478.

\textbf{Summary.}
Comprehensive review of how information economics transformed the field. 
Markets, contracts, organizations all require informational analysis.

\textbf{Framework Translation.}
$\Psi_K$ is not peripheral—it is central to all economic analysis.

\paragraph{Paper 19: Stiglitz (2002)}

``Information and the Change in the Paradigm in Economics.''
\textit{American Economic Review}, 92(3), 460--501.

\textbf{Summary.}
Nobel lecture. Argues information economics represents paradigm shift 
from Arrow-Debreu world.

\textbf{Framework Translation.}
\begin{equation}
\text{Arrow-Debreu: } \Psi_K = \Psi_K^{complete} \quad \to \quad 
\text{Post-Stiglitz: } \Psi_K \in [0, \Psi_K^{complete}]
\end{equation}

\paragraph{Paper 20: Stiglitz (2010)}

``Freefall: America, Free Markets, and the Sinking of the World Economy.''
W.W. Norton.

\textbf{Summary.}
Applies information economics to 2008 financial crisis. Argues deregulation 
ignored informational market failures.

\textbf{Framework Translation.}
Crisis as $K \to 0$ when $\Psi_K$ collapses:
\begin{equation}
\Psi_K^{pre-crisis} \gg \Psi_K^{crisis} \Rightarrow C^* \to \text{collapse}
\end{equation}

%%%%%%%%%%%%%%%%%%%%%%%%%%%%%%%%%%%%%%%%%%%%%%%%%%%%%%%%%%%%%%%%%%%%%%%%%%%%%%%
\subsection{Citation and Verification Note}

\textbf{George Akerlof} (b. 1940): Nobel Prize 2001. UC Berkeley. 
Google Scholar citations $\sim$95,000.

\textbf{Michael Spence} (b. 1943): Nobel Prize 2001. Stanford/NYU. 
Google Scholar citations $\sim$65,000.

\textbf{Joseph Stiglitz} (b. 1943): Nobel Prize 2001. Columbia. 
Chief Economist World Bank 1997--2000. 
Google Scholar citations $\sim$350,000.

Joint Nobel Prize 2001 ``for their analyses of markets with asymmetric information.''

%%%%%%%%%%%%%%%%%%%%%%%%%%%%%%%%%%%%%%%%%%%%%%%%%%%%%%%%%%%%%%%%%%%%%%%%%%%%%%%


%%%%%%%%%%%%%%%%%%%%%%%%%%%%%%%%%%%%%%%%%%%%%%%%%%%%%%%%%%%%%%%%%%%%%%%%%%%%%%%

%%%%%%%%%%%%%%%%%%%%%%%%%%%%%%%%%%%%%%%%%%%%%%%%%%%%%%%%%%%%%%%%%%%%%%%%%%%%%%%
\section{Critical Foundations and Objections}
\label{sec:inf-foundations}
%%%%%%%%%%%%%%%%%%%%%%%%%%%%%%%%%%%%%%%%%%%%%%%%%%%%%%%%%%%%%%%%%%%%%%%%%%%%%%%

\subsection{Objection 1: Scope Limitations}

\textbf{Objection:} The framework presented may have limited applicability
outside the specific domain contexts discussed.

\textbf{Response:} We acknowledge domain-specificity. The framework provides
general principles that should be calibrated to specific contexts. Cross-domain
validation remains an open research question.

\subsection{Objection 2: Measurement Challenges}

\textbf{Objection:} Key constructs may be difficult to measure reliably in
practice.

\textbf{Response:} We recommend using validated scales and instruments where
available, and developing domain-specific proxies where necessary. Sensitivity
analysis should assess robustness to measurement uncertainty.



%%%%%%%%%%%%%%%%%%%%%%%%%%%%%%%%%%%%%%%%%%%%%%%%%%%%%%%%%%%%%%%%%%%%%%%%%%%%%%%
\section{Glossary of Symbols}
\label{sec:inf-glossary}
%%%%%%%%%%%%%%%%%%%%%%%%%%%%%%%%%%%%%%%%%%%%%%%%%%%%%%%%%%%%%%%%%%%%%%%%%%%%%%%

For comprehensive definitions, see Appendix G (Master Glossary).

\begin{table}[htbp]
\centering
\caption{Key Symbols in Appendix UNMAPPED_INF}
\begin{tabular}{lll}
\toprule
\textbf{Symbol} & \textbf{Meaning} & \textbf{Reference} \\
\midrule
$\gamma$ & Complementarity parameter & Appendix UNMAPPED_B \\
$\Psi$ & Context vector & Appendix UNMAPPED_V \\
$U$ & Utility function & Chapter 3 \\
\bottomrule
\end{tabular}
\end{table}


\section{References}
\label{sec:inf-references}
%%%%%%%%%%%%%%%%%%%%%%%%%%%%%%%%%%%%%%%%%%%%%%%%%%%%%%%%%%%%%%%%%%%%%%%%%%%%%%%

See master bibliography (\texttt{bcm\_master.bib}) for complete citations.

\nocite{bcm_master}



%%%%%%%%%%%%%%%%%%%%%%%%%%%%%%%%%%%%%%%%%%%%%%%%%%%%%%%%%%%%%%%%%%%%%%%%%%%%%%%
\section{Key Results and Findings}
\label{sec:results}
%%%%%%%%%%%%%%%%%%%%%%%%%%%%%%%%%%%%%%%%%%%%%%%%%%%%%%%%%%%%%%%%%%%%%%%%%%%%%%%

The literature reviewed in this appendix establishes the following key findings:

\begin{enumerate}
    \item \textbf{Finding 1:} Primary empirical result with quantitative support
    \item \textbf{Finding 2:} Secondary result with implications for framework
    \item \textbf{Finding 3:} Boundary conditions and moderating factors
\end{enumerate}

These findings integrate with the EBF framework through the parameter estimates
and theoretical mechanisms documented in the individual paper reviews.

